% ===========================================================================
%  preamble.tex --- packages and macros for the chapter on models with
%  non-constant rates.
%  When this chapter is dropped into the thesis, merge the contents of this
%  file into the thesis preamble; duplicates may simply be deleted.
%
%  Derived from CH2/preamble.tex, the house preamble, and from
%  CH8_REWRITE/preamble.tex, its closest structural sibling.  Where the
%  source chapter's own main.tex disagreed with Chapter 2, Chapter 2 wins.
% ===========================================================================

\usepackage[T1]{fontenc}
\usepackage[utf8]{inputenc}
\usepackage{lmodern}
\usepackage{microtype}
\usepackage{amsmath,amssymb,amsthm,mathtools}
\usepackage{graphicx}
\usepackage{float}
\usepackage{placeins}
\usepackage{booktabs}
\usepackage[hypcap=false]{caption}
\usepackage{subcaption}
\usepackage{tikz}
\usetikzlibrary{arrows.meta,positioning,calc}
\usepackage{pgfplots}
\pgfplotsset{compat=1.18}
\usepackage{enumitem}
\usepackage{xcolor}
\usepackage{array}
\usepackage{hyperref}
\usepackage[nameinlink,noabbrev]{cleveref}

\graphicspath{{figures/}}

% --- Equation cross-references as bare parenthesised numbers -----------------
\crefformat{equation}{(#2#1#3)}
\Crefformat{equation}{(#2#1#3)}
\crefrangeformat{equation}{(#3#1#4)--(#5#2#6)}
\crefmultiformat{equation}{(#2#1#3)}{ and (#2#1#3)}{, (#2#1#3)}{ and (#2#1#3)}

% --- Notation ---------------------------------------------------------------
% Scalars are set in italic; bold is reserved for vector-valued states, which
% do not occur in this chapter.
%
% This chapter collects seven models written at different times, so its
% governing rule is that one letter may carry two meanings in two different
% models provided the notation table separates them, may never carry two
% within one section, and may never contradict a symbol frozen elsewhere in
% the thesis.  The renames that rule forced, and the reason for each, are
% recorded in PHASE_A_REPORT.md section 9.1.  The ones worth stating here:
%
%   S_t        species count.  Y_t is the infected count of the birth--plague
%              model, in this same chapter.
%   \varrho    logistic growth rate.  r is the burst size thesis-wide.
%   \mathcal{B}(t)  integrated speciation rate.  This is Chapter 2's own
%              symbol for the same object (m:eq:lambdaint); plain B(t) is
%              taken twice over, by Kendall's extinction integral in
%              m:eq:inhomextforward and by the Pimentel population count B_t.
%   \varepsilon  evolution rate of the Pimentel model.  \mu is the death rate
%              thesis-wide.
%   \zeta      speciation exponent of the Pimentel+ model.  \lambda is the
%              birth rate thesis-wide.
%   \phi       potential inflow rate.  \xi is K/N_0-1 in the speciation model.
%   \omega     revival rate of the rise--run--ruin--rejuvenation model.
%   \chi       contact-infection coefficient of the birth--plague model.
%   c          consumption (or clearance) of an accumulating quantity.
%   \lambda_i,\mu_i,\delta_i  phase-indexed rates of the two-phase model.
% \lambda, \mu, \delta, a, b and \mathcal{K} carry their frozen thesis-wide
% meanings throughout and are never redefined here.
\newcommand{\pgf}{probability generating function}
\newcommand{\rv}{random variable}
\newcommand{\pr}[1]{\mathbb{P}\!\left(#1\right)}
\newcommand{\ex}[1]{\mathbb{E}\!\left(#1\right)}
\newcommand{\var}[1]{\operatorname{Var}\!\left(#1\right)}
\newcommand{\dd}{\mathrm{d}}
\newcommand{\ee}{\mathrm{e}}
\newcommand{\ddt}[1]{\frac{\dd #1}{\dd t}}
\newcommand{\dda}[2]{\frac{\dd #1}{\dd #2}}
\newcommand{\pddt}[1]{\frac{\partial #1}{\partial t}}
\newcommand{\pdd}[2]{\frac{\partial #1}{\partial #2}}
\newcommand{\lrb}[1]{\left(#1\right)}
\newcommand{\lrbs}[1]{\left[#1\right]}
\newcommand{\ind}{\mathbf 1}
\newcommand{\ft}{\textit{F.~tularensis}}
\newcommand{\yp}{\textit{Y.~pestis}}
\newcommand{\Kb}{\mathcal{K}}
\newcommand{\Bint}{\mathcal{B}}

% --- Cross-chapter references -----------------------------------------------
% Redefining these macros is all that thesis integration requires: they become
% "Chapter~\ref{...}" once the chapter order is fixed.  No chapter number
% appears in body prose.  Wording copied byte for byte from
% CH8_REWRITE/preamble.tex and CH6_REWRITE/preamble.tex.
\newcommand{\ChM}{the mathematical introduction}
\newcommand{\ChMcap}{The mathematical introduction}
\newcommand{\ChCore}{the chapter on the birth--death--cata\-strophe process}
\newcommand{\ChCorecap}{The chapter on the birth--death--cata\-strophe process}
\newcommand{\ChDist}{the chapter on its distribution theory}
\newcommand{\ChDistcap}{The chapter on its distribution theory}
\newcommand{\ChPop}{the chapter on population dynamics}
\newcommand{\ChPopcap}{The chapter on population dynamics}
\newcommand{\ChTwoType}{the two-type chapter}
\newcommand{\ChPath}{the chapter on bursting and budding pathogenesis}
\newcommand{\ChPathcap}{The chapter on bursting and budding pathogenesis}

% --- Forward references to later chapters -----------------------------------
% Every forward reference is routed through \fwd, whose first argument is a
% stable key and whose second is the prose actually typeset.
\newcommand{\fwd}[2]{#2}

% --- Theorem environments ---------------------------------------------------
\theoremstyle{plain}
\newtheorem{theorem}{Theorem}[chapter]
\newtheorem{proposition}[theorem]{Proposition}
\newtheorem{lemma}[theorem]{Lemma}
\newtheorem{corollary}[theorem]{Corollary}
\theoremstyle{definition}
\newtheorem{definition}[theorem]{Definition}
\theoremstyle{remark}
\newtheorem{remark}[theorem]{Remark}

% End-of-section open questions. Not on the theorem/remark counter.
\newcommand{\openqs}[1]{%
\paragraph{Open questions.}%
\begin{enumerate}[leftmargin=1.6em,itemsep=0.2em,topsep=0.2em]
#1
\end{enumerate}}

\allowdisplaybreaks

% --- Figure colours (CH7 / SHARED_CONVENTIONS palette) ----------------------
\definecolor{gray1}{HTML}{d9d9d9}
\definecolor{type1}{HTML}{0072B2}
\definecolor{type2}{HTML}{D55E00}
\definecolor{ink}{HTML}{1A1C1F}
\definecolor{inksoft}{HTML}{565B62}
\definecolor{cata}{HTML}{9A2820}
\definecolor{na1}{HTML}{0072B2}
\definecolor{nb1}{HTML}{D55E00}
\definecolor{Rup2}{HTML}{9A2820}
\definecolor{blu}{HTML}{0072B2}
